\documentclass[11pt]{article}
\usepackage[a4paper,margin=1in]{geometry}
\usepackage{subfiles}
\usepackage{amsmath,amssymb,bm}
\usepackage{hyperref}
\usepackage{booktabs}
\usepackage{array}
\usepackage{mathtools}

\newcommand{\rs}{r_{\mathrm{s}}}

\title{Appendix 7 (to main theory): Perihelion Precession in the One--Metric Bi--Conformal ISPG Geometry}
\author{}
\date{}

\begin{document}
\let\phi\varphi
\maketitle

\section*{Purpose and scope}
This appendix computes the perihelion advance of a test body in the main theory one-metric bi--conformal exterior. At leading post-Newtonian order the result matches GR because the PPN parameters satisfy $\gamma=\beta=1$ (Appendix~4). We present (i) the standard PPN route and (ii) the direct bi--conformal weak-field expansion route, and give a numerical value for Mercury.

\section*{Non-negotiable consistency with the main theory (one-metric formulation)}
We work strictly within the main theory posture:
\begin{itemize}
\item matter is minimally coupled to the single physical metric $g_{\mu\nu}$;
\item the physical exterior geometry is bi--conformal, $ds^2=-e^{\phi}c^2dt^2+e^{-\phi}d\sigma^2$;
\item the exact static vacuum exterior scalar profile is $\phi(r)=-\rs/r$ with $\rs\equiv 2GM/c^2$ (the static exterior solution of the main theory~\cite{Lotishvili2026}).
\end{itemize}
No disformal mapping and no frame change are introduced.

\section{Standard perihelion advance in PPN form}\label{sec:ppn_app7}
For a nearly Keplerian orbit of semi-major axis $a$ and eccentricity $e$, the 1PN perihelion advance per orbit in the PPN framework is
\begin{equation}
\Delta\varpi_{\rm PPN}
=\frac{6\pi GM}{a(1-e^2)c^2}\left(\frac{2+2\gamma-\beta}{3}\right).
\label{eq:ppn_peri}
\end{equation}
In GR, $\gamma=\beta=1$ and the bracket equals unity.

Appendix~4 shows that the main theory bi--conformal weak-field geometry yields
\begin{equation}
\gamma=1,\qquad \beta=1.
\label{eq:gb}
\end{equation}
Inserting \eqref{eq:gb} into \eqref{eq:ppn_peri} gives the standard result
\begin{equation}
\boxed{\ \Delta\varpi=\frac{6\pi GM}{a(1-e^2)c^2}\ }\qquad (1{\rm PN}).
\label{eq:peri_final}
\end{equation}

\section{Direct check from the bi--conformal expansion}\label{sec:direct}
For completeness, we record the weak-field expansion of the main theory exterior:
\begin{align}
g_{tt}&=-e^{\phi}=-\exp\!\left(-\frac{\rs}{r}\right)
= -\left(1-\frac{\rs}{r}+\frac{\rs^2}{2r^2}+O(\rs^3)\right),\\
g_{rr}&=e^{-\phi}=\exp\!\left(\frac{\rs}{r}\right)
= 1+\frac{\rs}{r}+\frac{\rs^2}{2r^2}+O(\rs^3),
\end{align}
in isotropic spherical coordinates ($d\sigma^2=dr^2+r^2d\Omega^2$). At 1PN order, the effective Newtonian potential $U(r)=GM/r$ enters with $\gamma_{\rm PPN}=1$ (i.e.\ the spatial and temporal sectors carry equal contribution, matching the main-text notation $U_{\rm main} = GM/(c^2 r)$ via $U = c^2 U_{\rm main}$), which is precisely the condition that leads to \eqref{eq:peri_final}. The distinctive ISPG signature therefore appears at higher order (2PN), treated in the dedicated Shapiro and light-bending appendices.

\section{Numerical value for Mercury}\label{sec:mercury}
For Mercury, take $a=5.7909\times 10^{10}\,{\rm m}$, $e=0.2056$, and $GM_\odot=1.3271244\times 10^{20}\,{\rm m^3/s^2}$. Equation~\eqref{eq:peri_final} gives
\begin{equation}
\Delta\varpi_{\rm Mercury}\simeq 5.02\times 10^{-7}\ {\rm rad/orbit}
\simeq 0.1035^{\prime\prime}\ {\rm per\ orbit}
\simeq 42.98^{\prime\prime}\ {\rm per\ century},
\end{equation}
in agreement with the classic GR value.

\section*{Takeaway}
In the main theory one-metric bi--conformal formulation, the Solar-System PPN parameters satisfy $\gamma=\beta=1$ (Appendix~4). Consequently the perihelion advance matches GR at 1PN order:
\[
\Delta\varpi=\frac{6\pi GM}{a(1-e^2)c^2}.
\]
The ISPG discriminators are pushed to higher order and are provided by the dedicated 2PN Shapiro and 2PN bending appendices.

\end{document}
